SciGUI is a graphical user interface that can be used for Corvus. 
The installers (install.bat, install.sh) will install it automatically,
or you can get it from the github page \url{https://github.com/times-software/SciGUI/releases/latest}.
The GUI can be used to control corvus, open input files, plot output, and setup calculation
in an easy manner. If you installed with the install scripts, there will be a launcher
on your desktop called Corvus-version\_GUI.bat or Corvus-version\_GUI.command (windows, MacOS).
Once launched, you should see a main window with three panels. The left panel shows keyword buttons, which 
can be selected. The right screen shows help for the selected keyword. The top panel has some controls (run and plot) 
as well as filters for selecting which keywords to show on the left panel.

\subsection{General Use}
The GUI shows keywords on the left panel. When a user selects a keyword, 
the right panel shows help for that keyword, along with and "enable" checkbox,
and fillable fields for value input. If a user want to fill a keyword, they 
check the "enable" checkbox, then fill the fields. When the user is finished 
filling out all of the keywords needed for their calculation, they click the run
button. SciGUI will then write a corvus input file with all of the enabled keywords
filled out, then run corvus on the input file. Note that there are many possible
keywords, so there are filters to select keyword by subject or code, and also a search
bar to search for keywords. Finally, there is a "Quick Start" button to help users
get started with specific types of calculations.

\subsection{The controls and filter panel (top panel)}
This panel has the run button, the Quick Start selection, and a variety of ways to filter and search keywords:
\begin{itemize}
\item Filter by code
Which code is this a specific keyword for, or is it a general keyword.
\item Filter by subject
Keywords related to a specific property of subject.
\item Search bar
Search for specific terms within a keyword
\item "Only enabled" checkbox
Show only the enabled keywords.
\end{itemize}

\subsection{The Quick Start Menu}
If a user selects one of the options from the "Quick Start" menu, the 
target\_list keyword (and maybe some others) will automatically be enabled and filled out.
In addition, keywords will be colored according to importance:
\begin{itemize} 
\item RED: Required
\item BLUE: Most useful
\item GRAY: Other/Advanced
\end{itemize}
By default, when the "Quick Start" option is chosen, only required and useful cards are shown.
In order to fill in the minimim requirements for this type of calculation, the user then just 
needs to click on each of the red keywords shown in the keyword panel, and fill in the values
for that keyword. 

\subsection{Validation}
If a user has enabled a keyword (selecting it on the left panel, then clicking the "enable" checkbox on the right)
then they must fill all required fields or uncheck the enabled checkbox before selecting another keyword. Otherwise
a pop-up message will warn the user, and the same keyword will be selected again.
